SOAP adapts in the eigenbases of its matrix preconditioners and
refreshes them periodically, here every 10 optimizer ages by a
warm-started power step followed by QR (QR10). Calls that perform a QR
refresh account for about two-thirds of the reference optimizer's
measured time. We describe TurboSOAP, which keeps that refresh cadence
but makes four of every five refreshes approximate. A warm refresh
applies a basis-tracking step, a damped Jacobi-type rotation of the
current basis followed by two Newton--Schulz orthogonality corrections,
and every fiftieth age uses a QR repair (warm10/QR50). The tracking step
needs no decomposition; its matrix products take BF16 operands with
FP32 accumulation, whereas QR and eigendecomposition require FP32
inputs. In the same compiled BF16 implementation, a call with a warm
refresh costs 61.6 ms of optimizer time against 601.2 ms for a call
with a QR refresh. Our implementation also stores all persistent matrix
state in BF16 and compiles the matrix updates, retaining FP32 master
parameters and decomposition islands. On one 109-million-parameter
language model trained on a fixed OpenWebText cohort, two
model-initialization seeds finish 10,000 optimizer calls with
cross-entropy increases of 0.01124 and 0.01008 relative to a stock-beta
FP32 SOAP QR10 reference, corresponding to 1.13\% and 1.01\% higher
perplexity, after larger observed gaps of 0.118--0.124 CE at call 500.
Persistent matrix state falls from 2.320 to 1.160 GiB. Complete-step
time falls 18.66\% in a paired 9,900-call window for one seed and
17.36\% in a matched 100-call prefix for the other. That reference
differs in implementation as well as cadence. Against QR10 in the same
compiled BF16 implementation, timed in one launch, TurboSOAP's step
time is 12.1\% lower (320.5 against 364.8 ms). Exact refreshes help early: at
call 500 this matched QR10 is 0.098--0.116 CE ahead. At call 10,000 it
ends 0.0018 and 0.0003 CE above TurboSOAP on the two seeds, within
their 0.003 spread. Within the BF16 implementation, then, approximating
the refreshes costs no endpoint loss that two seeds can resolve,
whereas with FP32 state it costs 0.007 on the one seed measured. Policy
and implementation interact: moving to BF16 raises QR10's endpoint loss
by 0.010--0.013 but warm10/QR50's by 0.004, so the gap to the FP32
reference cannot be apportioned between them. This study does not
identify why. These measurements support a scoped
quality/time/memory tradeoff; they do not establish statistical
equivalence or attribute the speed improvement to any single
component.
